\subsection{The first coefficient beyond the exact bands}
\label{sec:first-boundary-normal-form}

At the next coefficient, the strict residual-character separation
used for the exact bands becomes an equality. We determine the
resulting necessary form of an unequal-character comparison.

\begin{proof}[Proof of Theorem~\ref{thm:first-boundary-normal-form}]
Write $k_0=(p^2+1)/2$, $b_0=(p-1)/2$, $e=k_0(p-1)p^r$ and
$n=r+3$. Since $p\nmid k_0$, one has $v_p(e)=r$. The inequality
$k_0<p^2$ gives
\[
 e<p^{r+3}=p^{2n-3-r}.
\]
Corollary~\ref{cor:all-index-two-slope-reduction} therefore reduces
an unequal-character comparison to two whole potentially
supersingular groups, both with two-slope good models.
Lemma~\ref{os-lem:minima} and Proposition~\ref{fc-thm:valuation}
give
\begin{equation}\label{os-eq:first-boundary-range}
 (p-1)p^{n-1}<z_i<pe=k_0(p-1)p^{n-2}.
\end{equation}
Both models are central: the lower bound divided by $e$ is
$p^2/k_0>1$. The upper bound is less than $p^{n+1}$, so
$v_p(z_i)\le n$.

\smallskip\noindent
\emph{A parameter of valuation $n$.}
Suppose both valuations were at most $n-1$, ordered $s\le t$.
Different leading tame eigenvalues would contradict
$s+t\ge2n$ in Lemma~\ref{os-lem:two-chains}. For equal leading
eigenvalues, write $z_1=p^sa$ and $z_2=p^tb$, with $p\nmid ab$.
The same lemma makes the nonzero integer $a-b$ divisible by
$(p-1)p^{2n-1-s-t}$, hence by $(p-1)p^{n-s}$. But
\[
 |a-b|<k_0(p-1)p^{n-2-s}<(p-1)p^{n-s}.
\]
This is impossible. Relabel the model of valuation $n$ as
\begin{equation}\label{os-eq:first-boundary-high}
 z_1=bp^n,\qquad
 1\le b<\frac{k_0(p-1)}{p^2}
       =b_0+\frac{b_0}{p^2}.
\end{equation}
Since $0<b_0/p^2<1$, it follows that $1\le b\le b_0$.
Every parameter of valuation $n$ has its coefficient in this
same interval. All valuations just used are less than $2n-1$,
as required in Lemma~\ref{os-lem:two-chains}.

\smallskip\noindent
\emph{Residual inertia characters.}
Let $m=p^2-1$, let $\zeta$ have order $m$, write
$\alpha_i=\zeta^{a_i}$, and put $\delta=\zeta^{p+1}$, of order
$p-1$. The leading eigenvalues and the canonical and quotient
residual inertia characters are
\begin{equation}\label{os-eq:first-boundary-residual}
 \begin{aligned}
 \lambda_i&=\alpha_i\zeta^{-pz_i}\in\mathbf F_p^{\times},\\
 \kappa_i&=\zeta^{(p+1)e-z_i-a_i},\qquad
 \eta_i=\zeta^{pz_i-a_i}=\lambda_i^{-1},\qquad
 \kappa_i/\eta_i=\delta^{-z_i}.
 \end{aligned}
\end{equation}
These are the arbitrary-index leading-term formulas explained in
\eqref{os-eq:half-square-residual}; the last equality uses
$(p-1)\mid e$. We use the formulas on inertia, without asserting
equality of any unramified character factors. The canonical lines,
and their images under the given torsion isomorphism, are actual
$K$-stable lines. The isomorphism identifies the unordered residual
composition-character pairs on inertia.

\smallskip\noindent
\emph{Equal leading eigenvalues cannot occur.}
Suppose $\lambda_1=\lambda_2$. If $v_p(z_2)=n$, write
$z_2=b'p^n$ with $1\le b'\le b_0$. Lemma~\ref{os-lem:two-chains}
gives $(p-1)\mid b-b'$. This interval forces $b=b'$, and equal
parameters and leading eigenvalues give equal tangent scalars and
special characters, a contradiction. Thus $t=v_p(z_2)\le n-1$.
The same lemma gives an integer $j$ such that
\begin{equation}\label{os-eq:first-boundary-partner}
 z_2=bp^t+j(p-1)p^{n-1}.
\end{equation}
The lower bound in \eqref{os-eq:first-boundary-range} implies
$j\ge1$, because $bp^t<(p-1)p^{n-1}$. Its upper bound and the
positivity of $bp^t$ imply
\[
 j<\frac{k_0}{p}<b_0+1,
 \qquad\text{and hence}\qquad 1\le j\le b_0<p.
\]

Equal leading eigenvalues give equal quotient inertia characters
$\eta_i$. Equality of the residual determinants gives equal
canonical inertia characters $\kappa_i$. Their ratio in the first
model is $\delta^{-b}\ne1$, since $1\le b\le b_0<p-1$.
There is at most one stable line with inertia character $\kappa_i$:
two distinct such lines would span the representation and make
both composition characters equal to $\kappa_i$.
Consequently the actual isomorphism identifies the canonical lines.

Take these two matched canonical quotients.
Lemma~\ref{lem:actual-quotients} gives an actual comparison at level
$p^N$, where $N=r+2$, and the $K$-isogenies preserve the unequal
special characters. Since the original models are central,
Proposition~\ref{rc-prop:quotients} gives central, two-slope
quotient parameters
\begin{equation}\label{os-eq:first-boundary-quotients}
 \begin{aligned}
 A&=p^rU,& U&=k_0(p^2-1)-bp^3,\\
 B&=p^rV-bp^t,& V&=k_0(p^2-1)-j(p-1)p^2.
 \end{aligned}
\end{equation}
In particular $e<A,B<pe$. Since $U\equiv V\equiv-k_0\pmod p$,
we have $v_p(A)=r$. If $t<r$, then $v_p(B)=t$; if $t>r$, then
$v_p(B)=r$; and if $t=r$ but $p\nmid k_0+b$, again $v_p(B)=r$.

In each of these cases the ordered quotient valuations $q\le s$
satisfy $q,s\le r$. Their sum is less than $2N$ and both are less
than $2N-1$, so Lemma~\ref{os-lem:two-chains} forces equal leading
eigenvalues. The nonzero difference of their unit parts is then
divisible by at least
\[
 (p-1)p^{2N-1-q-s}\ge(p-1)p^{r+3-q}.
\]
Its absolute value is smaller than
\[
 pe/p^q=k_0(p-1)p^{r+1-q},
\]
which is strictly less than that divisor because $k_0<p^2$.
This is impossible. The lower-level argument uses the explicit
valuation hypotheses of the two-chain lemma, without transferring
the original numerical bound to level $N$.

The only case left is $t=r$ and $p\mid k_0+b$.
Put $c_0=(p+1)/2$. Since $k_0\equiv c_0\pmod p$ and
$1\le b\le b_0$, the positive integer $c_0+b$ is at most $p$.
Therefore $b=b_0$ and $k_0+b_0=pc_0$. The first unit part in
\eqref{os-eq:first-boundary-quotients} is now
\[
 U=(p^3-1)/2,
\]
and the second parameter is
\[
 B=p^{r+1}W,\qquad W=p\bigl[k_0-j(p-1)\bigr]-c_0.
\]
Since $p\nmid c_0$, one has $v_p(B)=r+1$ exactly.
The quotient valuations $r,r+1$ are less than $2N-1$ and their
sum is less than $2N$. The two-chain lemma therefore forces equal
leading eigenvalues and requires $p^2\mid U-W$. But the exact
identity is
\[
 U-W=p(p-1)j.
\]
Its $p$-adic valuation is one, since $1\le j\le b_0<p$.
This contradiction excludes equal leading eigenvalues.

\smallskip\noindent
\emph{The remaining form and the actual residual splitting.}
Now $\lambda_1\ne\lambda_2$. Since $v_p(z_1)=n$ and
$v_p(z_2)\le n$, the two-chain inequality forces $v_p(z_2)=n$.
Write $z_2=b'p^n$ with $1\le b'\le b_0$. The quotient inertia
characters $\eta_1,\eta_2$ differ, so equality of the unordered
composition-character pairs exchanges their entries:
\[
 \eta_2=\kappa_1,\qquad \kappa_2=\eta_1.
\]
Equation~\eqref{os-eq:first-boundary-residual} gives
$(p-1)\mid z_1+z_2$, hence $(p-1)\mid b+b'$. Since
$0<b+b'\le2b_0=p-1$, this forces $b=b'=b_0$ and
$z_1=z_2=b_0p^n$.

The equal parameters give
$\alpha_2/\alpha_1=\lambda_2/\lambda_1\in\mathbf F_p^\times$.
Equality of the finite determinant factors also gives
$(\alpha_2/\alpha_1)^{p+1}=1$. Its order thus divides both $p-1$
and $p+1$, and hence divides two. The ratio is not one, because
the leading eigenvalues differ. Therefore $\alpha_2/\alpha_1=-1$.

Finally $\kappa_i/\eta_i=\delta^{-b_0}=-1$, so the two inertia
characters in each model are distinct. The actual isomorphism
sends the canonical line of the first representation to a
$K$-stable line in the second of inertia character $\eta_2$.
This line differs from the second canonical line and is therefore
a stable complement. Applying the inverse isomorphism to the
second canonical line supplies the complementary stable line in
the first representation. Both residual representations split
over $K$. The two distinct-inertia summands are unique, so the
given isomorphism exchanges canonical and complementary lines
in both directions. This proves all the asserted necessary
conditions.
\end{proof}
